\documentclass[11pt]{article}
\usepackage[T1]{fontenc}
\usepackage{lmodern,amsmath,amssymb,amsthm}
\usepackage[margin=1in]{geometry}
\usepackage{microtype}
\usepackage[colorlinks=true,linkcolor=blue,citecolor=blue,urlcolor=blue]{hyperref}
\newtheorem{theorem}{Theorem}
\newtheorem{lemma}[theorem]{Lemma}
\newcommand{\R}{\mathbb R}
\newcommand{\Z}{\mathbb Z}
\newcommand{\M}{\mathbf M}
\newcommand{\F}{\mathbb F}
\newcommand{\HH}{\mathcal H}
\newcommand{\LL}{\mathcal L}
\newcommand{\restr}{\mathbin{\vrule height 1.4ex depth 0pt width .07ex\vrule height .07ex depth 0pt width .65ex}}
\newcommand{\spt}{\operatorname{spt}}
\newcommand{\BV}{\operatorname{BV}}
\newcommand{\vol}{\mathrm{vol}}
\title{Integral optimal flat-norm decompositions in codimension one}
\author{Kaiwen Jin\\\small\href{mailto:yui.ui.siki@gmail.com}{yui.ui.siki@gmail.com}}
\date{13 October 2026}
\begin{document}
\maketitle
\begin{abstract}
The real flat norm decomposes a current into a residual and the boundary of
a filling, minimizing their total mass. We prove that every compactly
supported integral $d$-current in $\R^{d+1}$ has an integral optimal
decomposition, including when its boundary is nonzero. The filling is a
scalar function of bounded variation. Rounding its values with one common
shift preserves the mean filling mass and does not increase the mean residual
mass. The integer multiplicities of the input are essential to the latter
calculation. A compactness argument then gives an integral optimum.
\end{abstract}

\section{The flat-norm question}
Let $\mathcal N_k^c(\R^n)$ denote compactly supported real normal $k$-currents,
and let $\mathcal I_k^c(\R^n)$ denote compactly supported integral currents.
We use Euclidean comass to define current mass, denoted by $\M$.
For a compactly supported integral $d$-current $T$ in $\R^{d+1}$, set
\begin{equation}\label{eq:flat}
\F(T)=\inf\bigl\{\M(X)+\M(S):
T=X+\partial S,\ X\in\mathcal N_d^c(\R^{d+1}),\
S\in\mathcal N_{d+1}^c(\R^{d+1})\bigr\}.
\end{equation}
Integral currents describe oriented surfaces with integer multiplicities.
An integral optimum would retain that interpretation in the two parts of the
decomposition. Minimization over the integral class alone does not determine
whether its value equals the real infimum in \eqref{eq:flat}.

Ibrahim, Krishnamoorthy and Vixie \cite{IKV} prove this equality and integral
attainment for one-currents in the plane (Corollary 3.8). Their introduction
also records the case of codimension-one boundaries. Their general
codimension-one theorem (Theorem 3.7) uses a subdivision hypothesis to
transfer an integral simplicial optimum to the continuous problem. We give
a direct argument using scalar BV functions, which does not require that
hypothesis.

\begin{theorem}\label{thm:main}
For every integer $d\ge1$ and every compactly supported integral $d$-current
$T$ in $\R^{d+1}$, there exist
$X_I\in\mathcal I_d^c(\R^{d+1})$ and
$S_I\in\mathcal I_{d+1}^c(\R^{d+1})$ such that
\[
T=X_I+\partial S_I,\qquad
\F(T)=\M(X_I)+\M(S_I).
\]
\end{theorem}
Here $\partial$ is the current boundary, $\M$ is mass, and $\F$ is the real
infimum \eqref{eq:flat}. The input may satisfy $\partial T\ne0$.
The conclusion asserts existence of an integral optimum and does not require
every real optimum to be integral.

The argument rounds a top-dimensional filling rather than approximating the
input by a triangulation. Off the input surface, the BV coarea formula controls
the variation created by rounding. On the surface, both one-sided values must
use the same shift: their rounded difference lies between two adjacent
integers. Since the input multiplicity is an integer, the absolute residual
is affine on that interval. This gives the cost estimate needed for
Theorem~\ref{thm:main}.

\section{Currents and scalar BV functions}
Put $n=d+1$ and $\vol=dx_1\wedge\cdots\wedge dx_n$.
For $z\in C_c^\infty(\R^n;\R^n)$, write $\omega_z=\iota_z\vol$.
Every finite-mass $(n-1)$-current $A$ has a vector Radon measure $\mu_A$ with
\begin{equation}\label{eq:vector}
A(\omega_z)=-\int z\cdot d\mu_A,\qquad
\M(A)=|\mu_A|(\R^n).
\end{equation}
The sign is chosen for the identity below; the mass identity follows from
the isometry between Euclidean vectors and $(n-1)$-forms. For $u\in L^1(\R^n)$,
define the $n$-current $[u]$ by
$[u](f\vol)=\int uf\,dx$, $f\in C_c^\infty(\R^n)$.
If $u\in\BV(\R^n)$, integration by parts gives
\begin{equation}\label{eq:dictionary}
\mu_{\partial[u]}=Du,\qquad
\M([u])=\|u\|_1,\qquad
\M(\partial[u])=|Du|(\R^n).
\end{equation}
The derivative $Du$ is distributional; its total variation is the BV
seminorm. We always choose Borel representatives of BV functions.

\begin{lemma}\label{lem:top}
Every compactly supported normal $n$-current in $\R^n$ is $[u]$ for a
compactly supported $u\in\BV(\R^n)$.
\end{lemma}
\begin{proof}
The finite-mass measure representation gives a finite signed measure $\sigma$
on $n$-forms. Finite boundary mass means that the distributional gradient
$D\sigma$ is a finite vector measure, by testing against $\omega_z$.
For nonnegative smooth mollifiers $\rho_\varepsilon$, let
$u_\varepsilon=\sigma*\rho_\varepsilon$. Convolution gives
\[
\|u_\varepsilon\|_1\le|\sigma|(\R^n),\qquad
\int|\nabla u_\varepsilon|\,dx\le|D\sigma|(\R^n).
\]
For $\varepsilon<1$ the supports lie in one bounded set. BV compactness
\cite[slides 29--33]{Lahti} gives an $L^1$-convergent subsequence with limit
$u\in\BV(\R^n)$. On the other hand, mollification converges distributionally
to $\sigma$, so $\sigma=u\LL^n$. Testing outside $\spt\sigma$ shows that $u$
vanishes there. The measure representation used here and in
\eqref{eq:vector} is recalled in \cite[slides 15--17]{Lahti}.
\end{proof}

For the integral input, \eqref{eq:vector} takes the form
\begin{equation}\label{eq:input}
\mu_T=m\nu\,\HH^{n-1}\restr M,
\end{equation}
where $M$ is a Borel countably $(n-1)$-rectifiable carrier, $\nu$ is a
measurable unit normal, and $m\in L^1(\HH^{n-1}\restr M;\Z\setminus\{0\})$.
The sign of $m$ incorporates the chosen orientation.
In particular $\HH^{n-1}(M)\le\M(T)<\infty$.

We use two classical BV formulas, in the complete forms stated in
\cite[Theorems A.9 and A.18]{RT}. For $u\in\BV(\R^n)$, localized coarea gives
\begin{equation}\label{eq:coarea}
|Du|(B)=\int_\R |D\mathbf1_{\{u>t\}}|(B)\,dt
\end{equation}
for every Borel $B$. On an oriented countably $(n-1)$-rectifiable set $M$,
the one-sided approximate $L^1$ limits $u_M^+,u_M^-$ exist and are finite
at $\HH^{n-1}$-almost every point, and
\begin{equation}\label{eq:trace}
Du\restr M=(u_M^+-u_M^-)\nu\,\HH^{n-1}\restr M.
\end{equation}
The vector-valued source statement gives the scalar one by applying it to
$(u,0,\ldots,0)$ and taking the first component. For compactly supported $u$
and the carrier \eqref{eq:input}, the source formulas on a bounded open ball
apply on any ball containing the supports in its interior. Indicators of
nonzero superlevels are constant near that ball's boundary, so extending
them by the same constant introduces no boundary measure. Thus
\eqref{eq:coarea} also applies to the full-space superlevel perimeters used
below.

\section{Rounding without increasing the cost}
For $0<\theta<1$, let $q_\theta(a)=\lfloor a+\theta\rfloor$.
For every $a\in\R$, writing $a=j+r$ with $j\in\Z$ and $0\le r<1$ shows that
$q_\theta(a)$ takes the values $j,j+1$ with probabilities $1-r,r$. Hence
\begin{equation}\label{eq:scalar-volume}
\int_0^1q_\theta(a)\,d\theta=a,\qquad
\int_0^1|q_\theta(a)|\,d\theta=|a|.
\end{equation}
For $a,b\in\R$ and $m\in\Z$, we also have
\begin{equation}\label{eq:scalar-jump}
\int_0^1|m-q_\theta(a)+q_\theta(b)|\,d\theta=|m-a+b|.
\end{equation}
Indeed, if $a-b=j+r$, then
$q_\theta(a)-q_\theta(b)=j+\lfloor\{b+\theta\}+r\rfloor$, where braces
denote fractional part. The difference lies between the adjacent integers
$\lfloor a-b\rfloor$ and $\lceil a-b\rceil$, and its mean is $a-b$.
The function $s\mapsto|m-s|$ is affine on that interval because its corner
$m$ is an integer. This proves \eqref{eq:scalar-jump}, including all signs.

\begin{lemma}\label{lem:round}
Let $\mu=m\nu\HH^{n-1}\restr M$ be as in \eqref{eq:input}, and let
$u\in\BV(\R^n)$ vanish almost everywhere off a compact set $K$.
There is a compactly supported $v\in\BV(\R^n;\Z)$ such that
\begin{equation}\label{eq:cost}
|\mu-Dv|(\R^n)+\|v\|_1
\le |\mu-Du|(\R^n)+\|u\|_1.
\end{equation}
\end{lemma}
\begin{proof}
Take $u=0$ outside $K$ and set $v_\theta=q_\theta(u)$.
Then $v_\theta=0$ outside $K$ and $|v_\theta|\le|u|+\mathbf1_K$.
Tonelli's theorem and \eqref{eq:scalar-volume} give
\begin{equation}\label{eq:mean-volume}
\int_0^1\|v_\theta\|_1\,d\theta=\|u\|_1.
\end{equation}
For $k\in\Z$, put $E_{k,\theta}=\{u>k-\theta\}$ and define
\[
A_\theta(B)=\sum_{k\in\Z}|D\mathbf1_{E_{k,\theta}}|(B).
\]
Coarea and the partition of $\R$ by the intervals $(k-1,k)$ yield
\begin{equation}\label{eq:sample-coarea}
\int_0^1 A_\theta(B)\,d\theta=|Du|(B)
\end{equation}
for every Borel $B$. For almost every $\theta$, every sampled level has
finite perimeter and $A_\theta(\R^n)<\infty$. Fubini also gives
$u(x)+\theta\notin\Z$ for $\LL^n$-almost every $x$, for almost every $\theta$.
For these parameters the sums
\[
\sum_{k=-N}^N
\bigl(\mathbf1_{E_{k,\theta}}-\mathbf1_{\{0>k-\theta\}}\bigr)
\]
converge almost everywhere and in $L^1$ to $v_\theta$. Their absolute values
are bounded by $|v_\theta|$. Their derivative series converges in total
variation, since its total variation norms sum to $A_\theta(\R^n)$.
Distributional testing therefore shows that $v_\theta\in\BV$ and
\begin{equation}\label{eq:variation}
|Dv_\theta|\le A_\theta\quad\text{as measures}.
\end{equation}
This includes all diffuse variation of $u$.

Let $a(x)=u_M^+(x)$ and $b(x)=u_M^-(x)$ be the traces in
\eqref{eq:trace}. At each point where they exist, the parameters satisfying
$a(x)+\theta\in\Z$ or $b(x)+\theta\in\Z$ form a countable set.
Since the traces have measurable representatives and $\HH^{n-1}(M)<\infty$,
Fubini gives one full-measure set of parameters avoiding these equalities
at almost every point of $M$.
For such a parameter, the traces of $v_\theta$ are $q_\theta(a)$ and
$q_\theta(b)$. To justify this despite the discontinuity of the floor
function, if $c+\theta\notin\Z$ choose $\delta>0$ on which $q_\theta$ is
constant around $c$. For all $s\in\R$,
\[
|q_\theta(s)-q_\theta(c)|\le(1+\delta^{-1})|s-c|.
\]
Averaging this inequality over a shrinking half-ball proves the assertion
from the approximate $L^1$ trace of $u$, with no boundedness assumption.
Applying \eqref{eq:trace} to $v_\theta$ now gives
\[
(\mu-Dv_\theta)\restr M
 =\bigl(m-q_\theta(a)+q_\theta(b)\bigr)\nu\HH^{n-1}\restr M.
\]
Consequently \eqref{eq:scalar-jump} implies
\begin{equation}\label{eq:mean-on}
\int_0^1|\mu-Dv_\theta|(M)\,d\theta
=\int_M|m-a+b|\,d\HH^{n-1}=|\mu-Du|(M).
\end{equation}
On $M^c$ the input measure vanishes, so
\eqref{eq:variation} and \eqref{eq:sample-coarea} give
\[
\int_0^1|\mu-Dv_\theta|(M^c)\,d\theta\le|Du|(M^c).
\]
Together with \eqref{eq:mean-volume} and \eqref{eq:mean-on}, this bounds
the mean total cost by the right-hand side of \eqref{eq:cost}.
The cost is measurable: its variation term is the supremum over a countable
$C^1$-dense family of smooth vector tests of supremum norm at most one on a
fixed bounded neighbourhood of the supports. Each evaluation is measurable
by the joint measurability of $v_\theta$ and its uniform $L^1$ bound.
Thus at least one of the common good parameters has cost at most that mean.
Choosing it proves the lemma.
\end{proof}

For integer-valued $v\in\BV$ with compact support, $[v]$ has integer density
on $\R^n$, finite mass $\|v\|_1$, and finite boundary mass $|Dv|(\R^n)$.
The classical boundary rectifiability theorem therefore makes $[v]$ integral;
see \cite[Theorem 1.1]{DelNin} for an explicit statement and proof, and
\cite[Section 1.1]{IKV}. It follows that $T-\partial[v]$ is integral too,
and its boundary is $\partial T$. By \eqref{eq:dictionary},
Lemma~\ref{lem:round} therefore rounds any real admissible decomposition to
an integral one without increasing its cost.

\section{Attainment with compact support}
\begin{proof}[Proof of Theorem~\ref{thm:main}]
Choose $r_0>0$ with $\spt T\subset\operatorname{int}B_{r_0}$.
We first confine real competitors to one fixed ball. Take a smooth
$\rho:[0,\infty)\to[0,1]$ that is one up to $r_0$ and zero after $r_0+1$,
and let $g(r)=\int_0^r\rho(t)\,dt$.
The smooth map
\[
f(x)=g(|x|)\frac{x}{|x|}\quad(x\ne0),\qquad f(0)=0,
\]
is the identity on $B_{r_0}$ and has image in $B_{r_0+1}$.
Its radial and tangential derivative eigenvalues are $g'(r)$ and $g(r)/r$,
both in $[0,1]$. Thus $\|Df\|\le1$.
For compactly supported currents, define $f_\#$ by smooth pullback, using a
cutoff equal to one near the current and its boundary. Pullback commutes
with exterior differentiation there, so $\partial f_\#S=f_\#\partial S$.
The derivative bound and the mass definition give $\M(f_\#A)\le\M(A)$ in
every degree. Since $f_\#T=T$, this operation preserves the decomposition
and decreases both cost terms. Global properness of $f$ is unnecessary for
these compactly supported inputs.

The trivial competitor $(T,0)$ gives $0\le\F(T)\le\M(T)$.
Take a real minimizing sequence with costs at most $\F(T)+1/j$, confine it
by $f_\#$, and write its fillings as $[u_j]$ by Lemma~\ref{lem:top}.
Apply Lemma~\ref{lem:round} to choose integer-valued $v_j\in\BV$, supported
in the same fixed ball, and put $X_j=T-\partial[v_j]$.
Both currents are integral and
\[
\M(X_j)+\|v_j\|_1\le\F(T)+1/j.
\]
Their BV norms are bounded because
\[
\|v_j\|_1\le\F(T)+1,\qquad
|Dv_j|(\R^n)\le\M(T)+\M(X_j)\le\M(T)+\F(T)+1.
\]
BV compactness gives a subsequence converging locally in $L^1$ to
$v\in\BV(\R^n)$. The common support makes convergence global in $L^1$
and keeps $v$ supported in the closed ball. Pass to a further subsequence
with $\sum_j\|v_j-v\|_1<\infty$. Tonelli then gives pointwise convergence
almost everywhere, so $v$ is integer valued.

Set $S_I=[v]$ and $X_I=T-\partial[v]$. The preceding boundary rectifiability
argument makes both currents integral and compactly supported. Global $L^1$
convergence gives $\|v_j\|_1\to\|v\|_1$ and current convergence
$X_j\to X_I$. Mass is lower semicontinuous under current convergence by its
supremum definition. Hence
\[
\M(X_I)+\M(S_I)
\le\liminf_j\bigl(\M(X_j)+\|v_j\|_1\bigr)=\F(T).
\]
The reverse inequality holds because this pair is admissible for the real
infimum. The resulting equality proves the theorem, with
$\partial X_I=\partial T$ throughout.
\end{proof}

\paragraph{AI assistance.}
OpenAI Codex assisted with proof development, source checking and writing.
Independent mathematical review is requested.

\begin{thebibliography}{9}
\bibitem{IKV}
S. Ibrahim, B. Krishnamoorthy and K. R. Vixie,
\emph{Flat norm decomposition of integral currents},
Journal of Computational Geometry \textbf{7}(1) (2016), 285--307.
\href{https://doi.org/10.20382/jocg.v7i1a14}{doi:10.20382/jocg.v7i1a14}.
\bibitem{RT}
R. Rossi and M. Thomas,
\emph{From an adhesive to a brittle delamination model in thermo-visco-elasticity},
ESAIM: Control, Optimisation and Calculus of Variations \textbf{21}(1)
(2015), 1--59.
\href{https://doi.org/10.1051/cocv/2014015}{doi:10.1051/cocv/2014015}.
\bibitem{Lahti}
P. Lahti, \emph{Minicourse on BV functions},
University of Oxford, lecture slides, 16 June 2016.
\href{https://www.maths.ox.ac.uk/system/files/attachments/BV_minicourse_Slides_Version3_0.pdf}{Author's slides}.
\bibitem{DelNin}
G. Del Nin, \emph{Boundary rectifiability and compactness of integral currents
via $BV$ functions}, arXiv:2606.11185v1 (2026).
\href{https://arxiv.org/abs/2606.11185v1}{arXiv:2606.11185v1}.
\end{thebibliography}
\end{document}
